\begin{thebibliography}{99}
\bibitem{BBCL}M. Benaïm, I. Benjamini, J. Chen \& Y. Lima - “A generalized Pólya’s urn with graph based interactions”, Random Structures Algorithms 46 (2015) no. 4, p. 614-634 | DOI
\bibitem{Benaim}M. Benaïm - “Dynamics of stochastic approximation algorithms”, in Séminaire de Probabilités, XXXIII, Lect. Notes in Math., vol. 1709, Springer, Berlin, 1999, p. 1-68 | DOI
\bibitem{CH21}Y. Couzinié \& C. Hirsch - “Weakly reinforced Pólya urns on countable networks”, Electron. Comm. Probab. 26 (2021), article ID 35, 10 pages | DOI
\bibitem{CL}J. Chen \& C. Lucas - “A generalized Pólya’s urn with graph based interactions: convergence at linearity”, Electron. Comm. Probab. 19 (2014), article ID 67, 13 pages | DOI
\bibitem{Deneu90}J.-L. Deneubourg, S. Aron, S. Goss \& J.-M. Pasteels - “The self-organizing exploratory pattern of the argentine ant”, Journal of Insect Behavior 3 (1990) no. 2, p. 159-168 | DOI
\bibitem{book_ant}M. Dorigo \& T. Stützle - Ant colony optimization, MIT Press, 2004
\bibitem{EFR}D. Erhard, T. Franco \& G. Reis - “The directed edge reinforced random walk: The ant mill phenomenon”, 2019 | arXiv
\bibitem{ants_89}S. Goss, S. Aron, J.-L. Deneubourg \& J.-M. Pasteels - “Self-organized shortcuts in the argentine ant”, Naturwissenschaften 76 (1989) no. 12, p. 579-581 | DOI
\bibitem{Yassine}Y. Hamdi - “A model of ants’ search for food by reinforced walks on graphs” (2020), UG internship report, https://gitlab.com/yassine\_hamdi/a-model-of-ants..., École polytechnique and University of Bath
\bibitem{WARM_new}C. Hirsch, M. Holmes \& V. Kleptsyn - “Warm percolation on a regular tree in the strong reinforcement regime”, 2020 | arXiv
\bibitem{HJ04}B. M. Hambly \& J. Jordan - “A random hierarchical lattice: the series-parallel graph and its properties”, Adv. in Appl. Probab. 36 (2004) no. 3, p. 824-838 | DOI
\bibitem{HK+}M. Holmes \& V. Kleptsyn - “Infinite WARM graphs II: critical regime”
\bibitem{Janson04}S. Janson - “Functional limit theorems for multitype branching processes and generalized Pólya urns”, Stochastic Process. Appl. 110 (2004) no. 2, p. 177-245 | DOI
\bibitem{KMS}D. Kious, C. Mailler \& B. Schapira - “Finding geodesics on graphs using reinforcement learning”, 2020 | arXiv
\bibitem{LGR}L. C. Le Goff \& O. Raimond - “Vertex reinforced non-backtracking random walks: an example of path formation”, Electron. J. Probab. 23 (2018), article ID 39, 38 pages | DOI
\bibitem{Lima}Y. Lima - “Graph-based Pólya’s urn: completion of the linear case”, Stochastic Dyn. 16 (2016) no. 2, article ID 1660007, 13 pages | DOI
\bibitem{LL}G. F. Lawler \& V. Limic - Random walk: a modern introduction, Cambridge Studies in Advanced Math., vol. 123, Cambridge University Press, Cambridge, 2010 | DOI
\bibitem{LP}R. Lyons \& Y. Peres - Probability on trees and networks, Cambridge Series in Statistical and Probabilistic Math., vol. 42, Cambridge University Press, New York, 2016 | DOI
\bibitem{Markov}A. A. Markov - “Sur quelques formules limites du calcul des probabilités”, Izv. Ross. Akad. Nauk Ser. Mat. 11 (1917) no. 3, p. 177-186, (in Russian)
\bibitem{current_ants}Q. Ma, A. Tero, T. Nakagaki \& D. J. T. Sumpter - “Current-reinforced random walks for constructing transport networks”, J. R. Soc. Interface (2013), article ID 20120864 | DOI
\bibitem{Pemantle}R. Pemantle - “A survey of random processes with reinforcement”, Probab. Surv. 4 (2007), p. 1-79 | DOI
\bibitem{WARM}R. van der Hofstad, M. Holmes, A. Kuznetsov \& W. Ruszel - “Strongly reinforced Pólya urns with graph-based competition”, Ann. Appl. Probab. 26 (2016) no. 4, p. 2494-2539 | DOI
\end{thebibliography}
